%% =====================================================================
%% Actividad 16:
%% 
%% Proyecto "Analizador de Sentimientos para Reseñas de Productos"
%% 
%% Entregables:
%% Entregar una tabla que contenga tres columnas: en la primera listar mínimo ejemplos de herramientas de IA contemporáneas, la segunda columna será para especificar el tipo de IA a la que pertenece el agente dependiendo de su capacidad y/o funcionalidad y, la tercera columna será para argumentar dicha la ubicación, en base a su capacidad de autonomía, aprendizaje, razonamiento, y alcance de tareas.
%% En sesión virtual grupal revisar la clasificación de los agentes, discutiendo en grupo o en equipo los argumentos.
%% =====================================================================
\documentclass[11pt,letterpaper]{article}
% cambia a [sinindice] si es corta o [indice] para mostrar
\usepackage[sinindice]{tecnm}        

\materia{Inteligencia Artificial y su Ética}
\actividad{Tema 4.1. Procesamiento del lenguaje natural \\ Actividad 16 - Analizador de Sentimientos \\ para Reseñas de Productos}
\profesor{M. Lilia Ureña Lugo}
%\fecha{15 de octubre de 2026}    % opcional; por defecto es la fecha de hoy

\begin{document}
\portada
\indicegeneral

%% ------------------------- CONTENIDO ---------------------------------

\section{Introducción}

La Inteligencia Artificial se clasifica principalmente según su capacidad y su nivel de autonomía \parencite{iae_u1_diseno}. En la actualidad, los sistemas desarrollados pertenecen a la denominada IA Débil (o estrecha), diseñada para ejecutar tareas específicas en las que a menudo supera el desempeño humano en velocidad y precisión. En un ámbito teórico y futurista se encuentra la IA Fuerte (o AGI, por sus siglas en inglés), que busca replicar la capacidad intelectual humana para razonar, planificar y adaptarse a entornos complejos; así como la Súper-Inteligencia Artificial, un concepto hipotético donde las máquinas superarían ampliamente el intelecto humano en todos los campos.

Adicionalmente, la IA se categoriza según su funcionalidad \parencite{iae_u1_diseno} en sistemas reactivos, de memoria limitada, teoría de la mente y autoconscientes. En la Tabla 1 se presenta el análisis comparativo de 20 herramientas contemporáneas y perspectivas teóricas, argumentando su ubicación a partir de sus niveles de autonomía, aprendizaje, razonamiento y alcance de tareas \parencite{iae_u1_cuaderno}.



\section{Conclusión}

El análisis de las herramientas contemporáneas demuestra que la totalidad de las aplicaciones comerciales y operativas actuales se enmarcan dentro de la IA Débil con arquitecturas de IA Reactiva o de Memoria Limitada. Aunque sistemas avanzados han optimizado la velocidad y exactitud en sus dominios, continúan careciendo de autoconciencia, intencionalidad y comprensión contextual genuina.

Por su parte, categorías como la IA Fuerte (AGI) y la Súper-Inteligencia permanecen como construcciones teóricas y horizontes de exploración tecnológica  que por el momento solo funcionan de manera hipotética.

\referencias
\end{document}